\begin{abstract}
We investigate whether the Discrete Symplectic Cosmology (DSC) framework---in which cosmic evolution is modeled as a non-autonomous symplectic map on a Planck-scale lattice with a $1/\ln^2(t)$ adiabatic cooling law---can reproduce key statistical features of the cosmic microwave background (CMB) temperature anisotropies.
Using a St\"ormer--Verlet integrator on two- and three-dimensional lattices, we demonstrate that: (i) the symplectic evolution produces near-Gaussian one-point statistics (skewness $= +0.07 \pm 0.30$, excess kurtosis $= -0.28 \pm 0.24$ over 20 independent realizations), while a dissipative coupled-map-lattice baseline fails with kurtosis $= -1.60$; (ii) acoustic-like oscillation peaks emerge naturally in the power spectrum $D(k) = k^2 P(k)$, with positions alignable to a mock $\Lambda$CDM reference via three lattice parameters; (iii) the $1/\ln^2(t)$ cooling produces a finite acoustic horizon through asymptotic freeze-out of wave propagation; and (iv) Bayesian optimization recovers all three lattice parameters from power spectra in a twin experiment (MSE $= 8.9 \times 10^{-5}$).
Full-sky maps generated by projecting a $128^3$ lattice onto a Mollweide sphere exhibit multi-scale temperature fluctuations visually resembling Planck observations, with pixel distributions close to those of the Planck SMICA map.
These results establish the DSC lattice as a minimal computational framework capable of reproducing several CMB phenomenological features from symplectic dynamics alone, though the kurtosis sign differs from Planck and the acoustic peak spacing requires parameter fitting.
\end{abstract}
